\chapter{Unicode width and grapheme clusters}
\label{ch:width}

\begin{aside}
``How many columns does this character take?'' has no general
answer in Unicode. \maya{} ships a width table compiled from the
Unicode East Asian Width property and adjusts at runtime to match
what your terminal actually does.
\end{aside}

\section{The width problem}

A terminal cell is a fixed-pitch square. Most Latin characters
fill one cell. Some characters fill two: CJK ideographs, full-
width punctuation, most emoji. Some characters fill zero: combining
marks (diacritics layered onto a previous cell), zero-width
joiners, variation selectors.

If the renderer's idea of width does not match the terminal's,
columns drift. After one wrong-width character, the rest of the
line is offset by a column. Run a screen of mixed CJK and emoji on
a terminal with the wrong width table and the result is a smear.

\section{East Asian Width}

The Unicode East Asian Width property assigns each code point one
of six values:

\begin{itemize}[leftmargin=*]
  \item \textbf{F (Fullwidth).} 2 columns. CJK full-width Latin,
    full-width digits.
  \item \textbf{W (Wide).} 2 columns. CJK ideographs, most
    emoji.
  \item \textbf{H (Halfwidth).} 1 column. CJK half-width katakana.
  \item \textbf{Na (Narrow).} 1 column. ASCII Latin.
  \item \textbf{A (Ambiguous).} 1 in most contexts, 2 in CJK
    fonts. The classic trouble.
  \item \textbf{N (Neutral).} 1 column. Most Latin scripts.
\end{itemize}

\maya{}'s table maps each code point to 0, 1, or 2 columns. The
ambiguous class is the hard one.

\section{Ambiguous characters}

Greek letters, Cyrillic, the box-drawing characters, and many
mathematical symbols are Ambiguous. A user with a CJK locale and a
font tuned for it will see them as 2 columns. A user in a Latin
locale will see them as 1.

\maya{} picks 1 by default and exposes a setting to flip:

\begin{lstlisting}
maya::set_ambiguous_width(2);  // CJK-locale default
\end{lstlisting}

The setting is global because it must match the terminal, and a
terminal has one font.

\begin{warning}
Setting the wrong ambiguous width is the most common cause of
``my UI looks broken on \emph{their} machine''. Default to 1 and
let users opt in via a config flag.
\end{warning}

\section{Combining marks}

A combining mark (\code{U+0301} COMBINING ACUTE ACCENT, for
instance) attaches to the preceding base character. The pair
\code{e + U+0301} renders as ``e-acute'' in one cell. The combining
mark itself has width 0.

\maya{} tracks the previous base when iterating a string. A
combining mark appends to that cell's code point but does not
advance the column.

A canvas cell holds only one code point. To support combining
marks, \maya{} accumulates the base + marks into a small
side-table indexed by cell. The diff hashes both base and marks
when comparing.

\section{Grapheme clusters}

A grapheme cluster is a user-perceived character: a base plus
all its combining marks, plus zero-width joiners, plus emoji
sequences. Unicode TR\#29 specifies the boundary rules.

\begin{idea}
\textbf{Grapheme cluster.} The unit a user thinks of as ``a
character''. \emph{Not} the same as a code point. \emph{Not} the
same as a byte. The string ``e-acute'' may be one code point (NFC) or
two (NFD), but it is always one grapheme.
\end{idea}

\maya{} uses a compact state-machine implementation of the TR\#29
rules. Each input code point advances a state; the state
identifies whether the boundary is between this code point and
the next. On a boundary, the previous cluster is committed.

Emoji ZWJ sequences (the family emoji, for instance:
\code{man + ZWJ + woman + ZWJ + girl}) are one cluster.
Most terminals render them as one wide cell.

\section{Wide cells and the trailing slot}

When \maya{} paints a wide cell at $(x, y)$ it also writes a
sentinel into $(x+1, y)$ marking it as the trailing half. The
diff treats the pair as one unit; the painter refuses overlaps
into the trailing slot.

If layout asks for a width-1 element to overlap a wide cell, the
wide cell is blanked (replaced with two spaces) and the new cell
is written. The screen never shows half a wide cell.

\section{Emoji}

Emoji are wide (W) by default. Many have skin-tone modifiers,
gender modifiers, and ZWJ extensions. The runtime rule is:

\begin{itemize}[leftmargin=*]
  \item One base emoji + any modifiers + any ZWJ extensions =
    one cluster, two columns.
\end{itemize}

The combinatorial expansion is huge (the ``family'' emoji has
hundreds of valid forms) but they all collapse to two columns.
\maya{}'s cluster machine handles this without enumeration.

\section{Variation selectors}

\code{U+FE0F} forces emoji presentation; \code{U+FE0E} forces
text. So \code{U+2603} SNOWMAN can be:

\begin{itemize}[leftmargin=*]
  \item \code{U+2603 + U+FE0E}: text snowman, 1 column.
  \item \code{U+2603 + U+FE0F}: emoji snowman, 2 columns.
  \item \code{U+2603} alone: undefined; usually text.
\end{itemize}

\maya{} honours the variation selector when present and falls back
to the East Asian Width category otherwise.

\begin{funfact}
The snowman, \code{U+2603}, has been in Unicode since 1.1
(1993). It became an emoji in Unicode 6.0 (2010). Variation
selector support arrived later still. The little guy has been
through a lot.
\end{funfact}

\section{Normalisation}

NFC, NFD, NFKC, NFKD. \maya{} does not normalise input strings
--- it stores what you give it. But for comparison (in text-input
caret movement, for instance) you may want NFC.

\maya{} ships \code{nfc(std::u32string\_view)} as an opt-in. It
is not on the hot path because most application text is already
NFC (Web browsers normalise; most editors emit NFC).

\section{Bidirectional text}

Right-to-left scripts (Arabic, Hebrew) flow logically left-to-
right in memory but visually right-to-left. \maya{} does not do
BiDi at the moment --- it would require a per-line shaping pass
and a logical-to-visual mapping. If you need BiDi, consider
\code{libfribidi}.

A simpler workaround for occasional RTL strings: pre-reverse them
yourself before passing to \maya{}. The result is correct for
display but loses caret semantics.

\section{Width queries}

\code{maya::width(std::u32string\_view)} returns the column count
of a string. \code{maya::width(char32\_t)} returns 0, 1, or 2 for
a single code point.

These are \code{constexpr} since the width table is a
\code{constexpr} array. You can use them in static asserts:

\begin{lstlisting}
static_assert(maya::width(U"hello") == 5);
static_assert(maya::width(U"\u4e2d") == 2);  // U+4E2D
\end{lstlisting}

\section{The lookup table}

The width table is roughly 4 KB compressed. It is generated from
\code{EastAsianWidth.txt} (Unicode data file) at build time. The
generator outputs a sparse representation: ranges of code points
with their width.

\begin{lstlisting}
struct Range { char32_t lo, hi; uint8_t w; };
constexpr Range table[] = {
  {0x0000, 0x001F, 1},
  // ...
  {0x4E00, 0x9FFF, 2},   // CJK Unified Ideographs
  // ...
};
\end{lstlisting}

Lookup is a binary search: \code{O(log N)} where N is roughly
1000 ranges. Sub-microsecond.

\section{Recap}

\begin{recap}
\begin{itemize}[leftmargin=*]
  \item Width is per code point, sometimes per cluster.
  \item Ambiguous characters need a locale-aware choice.
  \item Combining marks are width 0; grapheme clusters are the
    user's ``character''.
  \item Variation selectors flip text vs emoji presentation.
  \item Wide cells reserve a trailing zero-width slot.
  \item \maya{}'s width table is generated, compressed, and
    \code{constexpr}.
\end{itemize}
\end{recap}

\section*{Workshop}

\begin{enumerate}[leftmargin=*]
  \item Write a string like ``hello (U+4E2D)(U+6587) plus a flag emoji''
    to a text element. Resize the terminal and confirm the
    columns stay aligned.
  \item Toggle ambiguous-width and observe Greek and box-
    drawing characters change column count.
  \item Build a grapheme-cluster iterator over a string and
    print each cluster on its own line.
\end{enumerate}
